\chapter{外部效度、可迁移性与数据融合}
\label{ch:transportability}

\begin{learninggoals}
读完本章，你应当能够：区分样本代表性、效应修饰与形式可迁移性；使用标准化、加权和选择图组织跨总体推断；审查重叠、测量协调、制度变化与系统反馈。
\end{learninggoals}

\section{外部效度、可迁移性与数据融合}

\subsection{目标人群、标准化与变量选择}

一项随机试验可以在入组样本中无偏估计治疗效应，却不能仅凭随机化说明结果适用于未入组者。\term{内部效度}关心样本内处理比较是否有效；\term{外部效度}关心结论能否适用于其他人群、地点、时间或干预版本。
文献中常区分\term{普遍化}（目标人群包含试验样本）与\term{迁移}（试验样本和目标人群可能来自不同总体），但术语使用并不完全统一。更重要的是明确两个域：源域 $S=1$ 提供何种实验或观测数据，目标域 $S=0$ 要回答哪个因果目标，它们在哪些机制上不同。
平均效应不是处理的固定属性。若治疗效果随年龄、基线风险、医疗资源或依从策略变化，样本构成改变就会改变ATE。外推因此不是在结果末尾加一句“需要更多研究”，而是一个可以建模的因果问题。

\paragraph{从样本效应到目标人群效应}

设试验中处理随机化，$S=1$ 表示进入试验，目标为整个目标人群中的 $E\{Y(a)\}$。若给定基线协变量 $L$ 后，潜在结果与入组状态可交换：
\[
Y(a)\indep S\mid L,
\]
并且试验与目标人群在相关 $L$ 上都有支持，则
\[
E\{Y(a)\}=\sum_l E(Y\mid A=a,S=1,L=l)P(L=l).
\]

右侧用试验估计每个 $L$ 层的处理结果，用目标人群的 $L$ 分布重新标准化。若目标只包含 $S=0$ 的非试验成员，则权重分布换成 $P(L\mid S=0)$。这与内部效应的混杂调整形式相似，但可交换对象不同：内部效度要求处理组可比，迁移要求源域与目标域在给定效应相关变量后可比。
选择概率加权提供等价实现。试验中较少代表目标人群的单位获得较大权重。与处理权重一样，接近零的试验参与概率产生极端权重；若目标人群包含试验从未覆盖的协变量组合，非参数迁移不可识别。

\paragraph{哪些变量需要用于迁移}

不是所有预测结果的变量都是\term{效应修饰因子}。一个变量可以强烈预测未治疗风险，却不改变绝对处理效应；也可能在风险比尺度上无修饰、在风险差尺度上有修饰。迁移变量的选择依赖效应尺度与结构。
理想集合应阻断选择 $S$ 与潜在结果之间导致效应差异的路径。纳入基线结局预测因子常有稳健性收益，但过度细分会破坏重叠。若一个未测变量同时影响入组和处理反应，单靠已测人口学标准化仍不足。
机制知识在这里特别重要。若药物通过某受体发挥作用，而受体分布在两域不同，受体或其可靠代理是迁移关键；若两域只在与效果无关的行政变量上不同，无需为表面代表性调整。样本“看起来像”目标人群不是充分标准。

\subsection{选择图、多源融合与时间变化}

Pearl与Bareinboim用\term{选择图}在因果图上增加选择节点，标记哪些机制可能跨域变化\parencite{pearlbareinboim2014}。若 $S_X\rightarrow X$，表示 $X$ 的生成机制在源域与目标域不同；没有选择节点指向某变量，则假定其条件机制稳定。

假设源域有 $X$ 的随机试验，目标域只有观测数据，且选择只影响中介前的协变量 $Z$。若在适当干预图中
\[
Y\indep S\mid Z,\doop(X),
\]
则 $Z$ 是对迁移\term{S-可容许}的：源域中按 $Z$ 分层的实验效应可用目标域 $P(z)$ 加权。

更复杂情形可用do-calculus与选择节点推导\term{迁移公式}。某些目标域机制虽然变化，却能由目标观测数据直接估计；稳定的干预关系则从源实验带入。可迁移性因此不是要求两域完全相同，而是要求差异的位置已知且可由可用数据桥接。

\paragraph{多源数据融合}

现实证据可能包括一个小型随机试验、大型电子病历、多个国家调查和机制实验。\term{数据融合}寻求组合各源的互补信息：试验提供处理赋值的内部效度，观察数据提供目标人群分布、长期结局或稀有亚组。

融合前必须统一：

\begin{itemize}
  \item 变量定义与测量尺度；
  \item 时间零点、随访和删失规则；
  \item 处理版本与依从策略；
  \item 每个数据源的选择机制；
  \item 哪些结构方程被假定跨源稳定。
\end{itemize}

同名变量不保证同一构念。不同医院的“重症”可能由不同编码规则产生；不同国家的教育年限代表不同课程质量。若测量机制跨域变化，直接把变量视为共同节点会把测量差异误作因果异质性。
融合也可能降低而非提高可信度。一个内部偏差严重的大数据集会以高精度支配结果；贝叶斯或加权整合若没有明确偏差模型，只是把不相容来源平均。应先写出每个来源为哪个目标提供何种识别信息，再决定组合。

\paragraph{时间迁移与政策改变}

从过去迁移到未来不仅是协变量再加权。诊断技术、背景治疗、行为响应和制度规则可能改变因果机制本身。COVID-19早期治疗效果无法仅按年龄结构加权到后期，因为病毒变体、免疫背景和标准护理同时变化。
政策还可能改变系统均衡。小规模补贴试验估计的是部分覆盖时的个体反应，全国推广会改变价格、劳动力供给和社会规范。SUTVA下的个体效应无法直接迁移到不同政策饱和度。应把覆盖率、网络或市场状态纳入处理定义，或使用结构经济模型与干扰设计。

\subsection{敏感性、数值标准化与测量协调}

迁移假设通常比样本内随机化更难验证。可以进行：

\begin{itemize}
  \item 比较源域和目标域的协变量支持与参与概率；
  \item 在源域内部留出子群，检验迁移方法能否恢复已知结果；
  \item 对未测效应修饰因子的分布和交互强度设定情景；
  \item 比较多个效应尺度和处理版本；
  \item 使用目标域中的负对照或短期可观测结局检查机制稳定性；
  \item 若缺乏重叠，报告目标子人群或部分识别界限。
\end{itemize}

外部效度不应以“样本是否人口学代表”二元判断。一个非代表性样本在测得关键效应修饰因子并有良好重叠时可以迁移；一个表面代表样本若处理、测量或机制与目标应用不同，仍可能不适用。

\paragraph{数值例子：同一分层效应，不同总体效应}

设试验中年轻人占80\%、老年人占20\%；治疗使年轻人风险下降2个百分点，使老年人下降10个百分点。试验ATE为
\[
0.8\times(-0.02)+0.2\times(-0.10)=-0.036.
\]
目标人群年轻人与老年人各占一半，则在分层效应可迁移时，目标ATE为
\[
0.5\times(-0.02)+0.5\times(-0.10)=-0.06.
\]
试验内部估计没有错，直接搬用 $-3.6$ 个百分点却回答了错误人群。

若年龄只是基线风险预测因子，而风险差在两层相同，重新加权不会改变风险差；在风险比尺度上却可能因基线风险不同出现另一表现。因此效应修饰和迁移必须指定尺度。一个变量是否“必须调整”不能脱离目标尺度宣布。
若目标人群包含试验中完全没有的90岁以上人群，公式要求的分层结果不存在。用回归平滑外推不是非参数识别，而是增加关于年龄---效应曲线的模型。报告应把重叠内标准化与重叠外模型外推分开。

\paragraph{测量协调的因果含义}

两个数据源都含“血压”，源域使用诊室单次测量，目标域使用家庭一周均值。二者误差结构和对白大衣效应的敏感度不同。若把它们当同一 $L$，选择差异可能其实是测量差异；若血压用于识别效应修饰，错误会破坏迁移可交换性。
协调可使用验证子样本中同时测两种指标，建立测量桥接模型；也可把潜在真实血压和两个测量节点显式画出。简单z分数标准化只统一单位与边际分布，不证明构念等价。
结局也需协调。一个系统用住院编码，另一个用临床判定；治疗还可能改变就医概率。外推时必须区分真实疾病机制是否稳定与结局观测机制是否稳定。选择图可以分别把选择节点指向潜在结局和测量节点。

\paragraph{迁移失败诊断与公平责任}

若加权后的源域能复现目标域中处理前或已知短期关系，却无法复现目标结局，可能有三种原因：遗漏效应修饰、处理或结局定义不一致、结果机制真实变化。仅继续增加人口学权重无法区分。
可使用“桥接结果”：选择在两域都观察、对机制变化敏感的中间结局，检查模型；使用多源实验比较同一分层效应；或在目标域开展小型校准试验。若少量目标实验与大源试验不一致，可用分层模型估计变化，而非强行假定完全相同。
运输公式的失败也是信息：它指出哪个机制需要目标域数据。最有价值的新样本不一定扩大总量，而可能专门覆盖缺失协变量层或直接测量疑似变化机制。

\paragraph{公平与外部效度}

训练与试验样本往往对边缘群体覆盖不足。总体平均效果良好可能掩盖某群体没有重叠或测量质量更差。以“总体可迁移”为唯一验收，会把不确定性集中到最缺乏证据的人。
报告应按政策相关群体展示支持范围、分层效应与迁移权重，不把极端权重产生的高方差估计包装为同等证据。必要时限制部署、优先收集目标群体数据，或采用保守决策界限。外部效度因此也有分配伦理：谁被纳入知识基础，会影响谁承担模型失败。

\section{从迁移公式到政策责任}

\subsection{迁移目标与选择图表达}

外部效度不是问研究“能否推广”这一笼统问题，而是定义目标总体中的效应。例如试验样本以$S=1$表示，目标总体以$S=0$表示，要估计
\[
E\{Y(1)-Y(0)\mid S=0\}.
\]
若给定处理前变量$X$后，试验参与与潜在结果独立，并且目标总体中的每类$X$在试验中都有支持，则可用试验内条件效应按目标分布标准化。随机化解决试验内部处理比较，参与交换性和参与正值性则连接试验与目标，两者缺一不可。

需要迁移的变量是效应修饰和选择机制的交集，而非所有结局预测因子。一个强预测因子若不改变处理效应，试验与目标分布不同不会造成平均效应差；一个弱边际预测因子却可能强烈修饰效应。领域知识、分层试验结果和机制证据共同决定变量选择，纯粹用预测重要性排序会偏离目标。
目标总体也可能随时间变化。把十年前试验迁移到今天，不只是年龄和疾病分布变化，治疗标准、诊断阈值、竞争风险和实施能力都可能改变。时间应作为环境机制的一部分，而非仅在回归中加入年份。

\paragraph{选择图怎样组织跨域差异}

选择图用特殊节点指向在源域和目标域之间可能改变的机制。若选择节点只指向基线变量$X$，而$Y$的结构机制给定$A,X$后稳定，可以通过目标域$P(X)$与源域试验条件效应组合。若选择直接指向$Y$，说明结果机制存在未由测量变量解释的变化，简单重加权通常不足。
图的价值是区分分布变化与机制变化。人口年龄结构改变属于输入分布变化；同样治疗在新医院因执行质量不同而改变效果，属于处理或结果机制变化；结局编码系统改变则是测量机制变化。三者可能产生相似的观测差异，却需要不同修复：重加权、重新校准实施模型或做测量桥接研究。
多源数据融合可让不同数据源贡献不同模块。随机试验提供处理到结果的内部效应，登记数据提供目标协变量分布，机制研究约束中介路径，政策数据库提供环境差异。但来源之间必须有共同变量语义和可连接单位；若“严重度”在各库中定义不同，形式上的同名节点不能直接拼接。

\subsection{测量协调与规模效应}

跨域研究常把测量差异当作数据清洗问题。实际上测量过程可能决定因果目标。试验中的抑郁结局由临床访谈获得，目标系统只有处方记录；二者不仅噪声大小不同，还可能测量不同构念。若处理影响求医和处方机会，处方记录甚至是处理后选择变量。
协调应建立构念---指标图：先定义希望比较的潜在构念，再记录每个数据源的观测指标、阈值、测量时点和缺失机制。可以使用共同样本、校准研究、潜变量模型或多指标设计建立桥梁，但每种方法都增加假设。没有桥接数据时，把变量重命名为同一英文缩写不会创造可比性。
处理也需要协调。一个试验的“常规照护”可能包含频繁随访，目标地区的常规照护几乎没有服务；药物剂量、供应、依从支持和共同治疗都会改变版本。迁移结论应针对完整策略，并说明哪些组成部分被假定等价。

\paragraph{规模扩大为何可能改变效应本身}

小规模试验通常把资源价格、人员行为和网络环境视为背景不变。全面部署会改变这些背景：培训需求降低执行质量，补贴推高价格，推荐系统改变内容供给，疫苗覆盖产生群体免疫。个体潜在结果因此依赖总体处理比例或网络配置，简单把试验平均效应乘以人口规模会失败。
处理这种问题需要定义饱和度或政策水平。可比较不同集群覆盖率下的直接效应与溢出效应，或建立供需和网络传播模型。外推对象不再是“同一个体处理从0变1”，而是整个系统从一种配置转向另一种配置。识别需要集群随机化、分阶段实施或结构模型，并清楚说明均衡假设。
规模效应也可能改变受益者组成。试点期最有动机、最易服务的人先参与，扩展后进入边缘群体；早期执行者经验丰富，后期机构能力较弱。监测部署过程应持续更新目标总体、实施版本和支持区域，而不是把一次外部效度分析视为永久许可证。

\paragraph{外部效度的分配责任}

哪些人进入试验、哪些结局被测量、哪些亚组有足够支持，会决定谁能从证据中受益。若弱势群体长期被排除，模型对他们的不确定性更高；按预测置信度分配资源又可能进一步排除他们，形成证据贫困循环。外部效度因此不仅是技术准确性，也是研究资源的分配问题。
负责任的报告应展示目标群体覆盖、重叠和分层不确定性，不用总体平均掩盖支持缺口。部署可以采用分阶段策略：在证据充分群体中先实施，在缺口群体中同步开展有保护措施的研究；也可以使用保守决策界限，避免把高方差外推包装成个性化精确。
价值判断还进入目标选择。最大化总体平均收益可能牺牲小群体，最小化最坏群体风险则可能降低平均效率。统计模型不能独自选择公平准则，但可以明确不同准则需要哪些效应、成本和不确定性信息，使规范分歧不再隐藏在默认损失函数中。

\section{迁移分析的计算、诊断与治理}

\subsection{标准化计算与普遍化诊断}

设试验中低风险层占80\%，处理效应为风险下降2个百分点；高风险层占20\%，处理效应为下降10个百分点。试验平均效应是$0.8\times(-0.02)+0.2\times(-0.10)=-0.036$。目标人群若低风险占30\%、高风险占70\%，标准化效应为$0.3\times(-0.02)+0.7\times(-0.10)=-0.076$。
两个数值都正确，却回答不同总体。差异不是试验“失去内部效度”，而是效应修饰与人群组成共同作用。若只报告相对风险，而基线风险也跨域变化，绝对收益差异可能更大。迁移报告应同时展示分层效应、源域与目标分布以及加权后的总体结果。
如果目标高风险层在试验中完全没有样本，上式中的$-0.10$无法由试验估计。用回归预测出的数值依赖外推模型；应限制目标、补充试验或给出基于机制与外部数据的范围。正值性诊断比权重公式更先决定结论可信度。

\paragraph{试验普遍化的诊断顺序}

第一步比较资格标准：哪些目标成员原则上不能进入试验，排除是否与效应修饰相关。第二步比较参与倾向与协变量分布，展示试验和目标中的重叠。第三步检查处理、结局和随访是否同义。第四步估计参与权重或标准化模型，并报告权重分位数与有效样本量。第五步使用替代变量集、截尾和结果模型检查敏感性。
诊断不能只看加权后基线均值。参与模型可能在边缘群体外推，结局机制可能跨机构变化，未测效应修饰仍存在。可利用目标域观测结局、不同试验或负对照环境检查模型，但这些数据也有自身混杂。最可靠的结论会明确哪些差异已由数据桥接，哪些仍由稳定性假设承担。
普遍化与可迁移性有时区分为目标总体是否包含试验样本，但实质逻辑相近：都是把源域条件关系与目标域分布组合。术语选择不应遮蔽共同的交换与支持条件。

\subsection{证据冲突与部署后审计}

两个试验给出不同效果，可能因人群、剂量、对照、结局或执行质量不同，也可能只是抽样波动。数据融合前应建立差异表，而不是立即做一个加权平均。若差异可由测量的效应修饰解释，可分别标准化到共同目标；若处理版本不同，应保留为不同策略；若结果机制直接变化，层次模型的“研究随机效应”只能描述异质，不能解释可迁移性。
观察研究与试验冲突时，试验在内部处理比较上通常更强，但观察研究可能覆盖长期、罕见结局和不同实施环境。应检查试验依从、失访和结局时窗，同时对观察研究做偏差分析。冲突可生成新的目标：短期生物效应是否稳定，而长期政策效应为何随系统改变。
证据综合的输出可以是分层结论，而非单一数值。例如对试验覆盖人群给出较窄效应，对外推人群给出宽范围，对新实施版本暂不量化。这样的结构比一个忽略来源差异的总体元分析更适合决策。

\paragraph{部署后的持续迁移审计}

迁移不是发布前一次性步骤。部署会改变参与者、执行者和环境，原有目标分布不断更新。系统应监测处理版本、关键效应修饰、重叠、结局测量和策略依从；当这些指标越过预定阈值，触发重新估计、局部试验或暂停外推。
结果监测要避免选择反馈。政策决定谁接受处理，接受又决定谁产生标签；只比较部署后处理组与未处理组会重现混杂。可保留小比例随机探索、使用分阶段推广或设计可信的准实验。若伦理上不能随机，应至少保存完整决策日志和处理概率，为离策略评价提供依据。
模型退役也是迁移治理的一部分。新治疗替代旧对照、编码体系改变或关键群体缺乏支持时，应明确结束旧结论的适用期，而不是无限再校准。证据有生命周期，外部效度报告应包含下一次审查条件。

\subsection{完整迁移案例：从专科试验到城市系统}

某专科中心完成了一项康复项目随机试验，现在城市医疗系统想判断全面推广后的一年住院风险差。$S=1$表示进入试验，$S=0$表示只在目标总体登记；$A$为康复项目，$Y$为住院，$L$包含基线风险与行动能力。试验保证$S=1$中$A$的可交换性，却不保证试验参与者与城市患者拥有相同结局机制。

\begin{center}
\begin{tikzpicture}[node distance=13mm and 17mm]
\node[causalnode] (s) {$S$};
\node[causalnode,right=of s] (l) {$L$};
\node[causalnode,right=of l] (a) {$A$};
\node[causalnode,right=of a] (y) {$Y$};
\node[causalnode,below=of a] (v) {$V$};
\draw[causalarrow] (s)--(l);
\draw[causalarrow] (l)--(a);\draw[causalarrow] (l) to[bend left=22] (y);
\draw[causalarrow] (a)--(y);
\draw[causalarrow,dashed] (s)--(v);\draw[causalarrow,dashed] (v)--(y);
\end{tikzpicture}
\end{center}

实线部分表示，试验招募使$L$的分布不同，而$L$修饰处理效应。若给定$L$后潜在结果可在试验与目标间交换，且目标每个$L$层都在试验有支持，则
\[
E{Y(a)\mid S=0}=
\sum_l E(Y\mid A=a,L=l,S=1)P(l\mid S=0).
\]
虚线部分表示一个更强风险：实施容量$V$在专科中心与城市系统中不同，并直接影响$Y$。若$V$没有测量，只对$L$标准化无法修复机制变化。这不是权重模型的小误差，而是小规模试验中的“每人每周两次指导”与全市推广时的处理版本已经不同。

\begin{workedexample}[总体组成改变的数值后果]
将$L$分为可独立行走和行动受限两层。试验分别估计的住院风险差为$-0.08$与$-0.18$。试验中两层占比为0.75和0.25，因而试验平均为
\[
0.75\times(-0.08)+0.25\times(-0.18)=-0.105.
\]
城市目标中两层占比为0.35和0.65，若分层效应可迁移，目标平均为
\[
0.35\times(-0.08)+0.65\times(-0.18)=-0.145.
\]
差异来自效应修饰与总体组成，不是试验随机化失效。但若城市中还有“需他人协助行走”而试验完全排除的第三层，该层效应不能由两个已观察数字插值。必须限定目标、补充试验或给出外推范围。
\end{workedexample}

\paragraph{参与权重的解释与诊断}

当试验参与可由$L$解释时，可为试验中的人赋予与其在目标中代表性成比例的权重。一种参与几率权重与
\[
\frac{P(S=0\mid L)}{P(S=1\mid L)}
\]
成比例。它使试验中稀少、目标中常见的行动受限者获得更大代表性。权重并不证明这些人与未参加者在所有效应修饰上相同；它只实施已写出的可迁移假设。

诊断应先分层回答三个问题。第一，目标中有哪些人几乎不可能进入试验？第二，大权重是因为少数合理的稀有类型，还是参与模型函数形式不稳定？第三，加权后已测$L$是否与目标分布对齐？应报告权重分位数、有效样本量、每个目标层的试验人数和截尾后目标的变化。

\paragraph{测量协调先于数值融合}

专科中心可能以专业评定表测量行动能力，城市登记系统却只有辅具申请和诊断编码。即使两边都有一个名为“行动限制”的二元变量，切点、记录动机和测量时点不同也会使标准化失真。应先建立变量映射表，报告可互换、需校准和无法对齐的字段；再用重叠子样本或验证研究估计测量桥接。

结局测量也需同样审查。试验中的住院由独立委员会裁定，目标数据中则来自医保结算；跨市就医或康复项目对就医行为的改变会造成差异漏报。若测量误差受$A$或$S$影响，它不是简单降低精度，而会改变迁移查询的含义。

\begin{center}
\small
\begin{tabularx}{\textwidth}{>{\bfseries}p{2.1cm}X X}
\toprule
风险层 & 可观测诊断 & 应对动作\\
\midrule
总体支持 & 目标层在试验无或极少样本 & 限定目标、补充招募或报告外推范围\\
效应修饰 & 分层效应差异且人群组成改变 & 按目标分布标准化并报告分层结果\\
测量协调 & 字段同名但时点、切点或来源不同 & 建立映射与验证子研究\\
实施版本 & 容量、依从或对照治疗跨域变化 & 改写处理策略并进行局部试验\\
系统反馈 & 覆盖率扩大后等待、价格或行为变化 & 监测饱和度并使用分阶段部署\\
\bottomrule
\end{tabularx}
\end{center}

\begin{failurecheck}[把内部效度高写成全部人群可用]
随机化使试验内的处理比较可信，但不回答目标中缺失的效应修饰、处理版本或规模效应。“高质量RCT已证明有效”至少应拆成三句：试验内哪个策略比较被识别，目标总体与试验差异由哪些变量桥接，全面部署是否保留相同干预与结局机制。任何一句缺少时，数值都应带有相应的限定。
\end{failurecheck}

\subsection{多源数据融合：每个来源只提供它识别的部分}

现实中常没有一份数据同时包含随机处理、长期结局、代表性人群和完整测量。数据融合的任务不是把所有行连接成一张大表，而是明确每个来源识别哪个因子或桥接关系。例如，随机试验可识别短期处理效应，城市登记可提供目标协变量分布与长期结局，一个小型验证子研究可连接两种测量标尺。

\begin{center}
\small
\begin{tabularx}{\textwidth}{>{\bfseries}p{2.2cm}X X}
\toprule
数据来源 & 主要可提供的信息 & 不能单独提供的信息\\
\midrule
专科中心RCT & 已定义处理版本下的内部处理比较 & 全市人群组成、大规模实施效应\\
城市登记 & 目标总体分布、常规实施与长期结局 & 无未测混杂的处理效应\\
测量验证子研究 & 专业评定与登记代理之间的误差桥接 & 处理效应与系统规模反馈\\
部署试点 & 目标环境中的依从、容量、安全与短期效果 & 未试点地区的无条件长期效应\\
\bottomrule
\end{tabularx}
\end{center}

一个可识别的融合式必须对这些信息接口有明确要求。例如要用RCT的短期行动改善$M$与登记中$M$到长期住院$Y$的关系推断$A$到$Y$，需要$M$足以中介处理的长期影响，登记中$M-Y$关系可在测量的共同原因后交换，且试验与登记中$M$的测量可比。任何一项都比“两个数据集样本大”更重要。

\paragraph{证据冲突要先对齐估计对象}

假设RCT报告风险差$-0.10$，登记调整研究为$-0.04$，部署试点为$-0.07$。不能先假定三个数字在估计同一个常数真值。应分别列出资格、处理版本、对照水平、随访时窗、结局尺度与效应人群。若RCT为密集专业康复，登记的处理为“至少一次转诊”，数值差异首先反映处理不同，而不是某一研究“错了”。

对齐后仍有冲突时，可把每个偏差方向显式化。专科试验依从较高，可能高估全面实施效果；登记中未测病情使重症者更容易接受康复，可能将保护效应向零拉；试点规模小，容量反馈尚未完全出现。这些方向不会自动告诉真值，但会指出需要的负对照、依从调查、容量实验与未测混杂敏感性参数。

\paragraph{规模改变会让干预本身内生变化}

在500人试验中，每位治疗师负责10人；全市推广后若仍只有50位治疗师，每人负责数上升、等待延长，项目实际剂量与质量会变化。此时不能说“个体效应不变，只是更多人参加”。覆盖率通过系统资源反过来改变处理版本。

可把治疗师负荷和等待时间放入干预描述，定义不同覆盖率下的总体效应；也可在不同地区分阶段提高覆盖率，直接学习饱和度反应。政策容量是因果变量，不是预算附注。对存在网络、市场与供给约束的系统，只估计低覆盖个体效应不足以支持全面部署。

\begin{failurecheck}[把数据源标签当作一个普通协变量]
在联合数据中加入“来源=RCT/登记”指示变量并拟合一个大回归，不会自动使两个数据生成机制可交换。来源节点可能同时改变招募、处理分配、结局测量与实施强度。应使用选择图或明确的数据融合假设分别标记哪些模块改变，再检查手中数据是否为每个改变提供桥接。
\end{failurecheck}

\begin{chapterreview}
内部效度不能自动保证外部效度。迁移把源域实验关系与目标域协变量分布或观测机制结合；选择图明确哪些机制跨域变化。可靠融合要求统一变量、时间和处理版本，并说明每个数据源提供的识别信息。效应修饰、重叠、测量变化和系统均衡是外推的核心限制。
\end{chapterreview}

\begin{selfcheck}
\begin{enumerate}
  \item 推导试验效应按目标人群协变量分布标准化的公式，并列出三项假设。
  \item 为什么强结局预测因子不一定是迁移所需的效应修饰变量？
  \item 为两个医疗系统画选择图，标出治疗、测量和基线风险的可能变化机制。
  \item 解释小规模政策试验推广到全国时，个体效应为何可能失效。
\end{enumerate}
\end{selfcheck}

\chapterreadings{
选择图与形式可迁移性见\textcite{pearlbareinboim2014}和\textcite{bareinboimpearl2016}。试验普遍化与加权方法可读\textcite{cole2010}、\textcite{stuart2011}和\textcite{westreich2017}；嵌套与非嵌套试验设计及其抽样含义见\textcite{dahabreh2019}与\textcite{dahabreh2021}。\textcite{degtiarrose2023}提供了外部效度方法的系统综述。将可迁移性置于目标试验框架中时，可结合\textcite{hernanrobins2020}。

形式可迁移性的一般判定算法见\textcite{bareinboimpearl2013}。基于倾向分层的试验普遍化见\textcite{tipton2013}，样本与总体相似度的量化诊断见\textcite{tipton2014}；不同推广估计器的模拟和经验比较可读\textcite{kern2016}。抽样权重在实际随机试验中的实现见\textcite{buchanan2018}，而\textcite{westreich2019}强调内部效度、目标效度和研究设计等级不能混为一谈。

跨地点预测政策效果的早期计量方案见\textcite{hotz2005}，教育干预普遍化方法的专题综述见\textcite{tiptonolsen2018}。这些文献应与正文的处理版本、重叠和规模反馈共同阅读，避免把统计重加权误当成机制自动稳定。

迁移不仅是抽样权重问题。从证据进入新政策环境所需的支撑因素见\textcite{cartwright2007}，随机试验为何不能仅凭设计标签保证外部适用性可对照\textcite{deatoncartwright2018}。
}
